\chapter{\IfLanguageName{dutch}{Probleem Stelling}{Problem}}
\label{ch:probleem_stelling}

\section{Huidige situatie}
Het Emissie Model voor Ammoniak in Vlaanderen (EMAV).
Een programma dat op basis van gegevens van landbouwers wiskundige berekeningen uitvoert en de ammoniak emmissies berekend.
De Vlaamse Milieumaatschappij (VMM) gebruikt het model als hun standaard rapporterings en voorspellings applicatie.
\autocite{DemeyerPEMAV}

Wat de vlaamse overheid exact wilt doen is beschreven door de VMM en de Vlaamse Landmaatschappij (VLM) in dit protocol \textcite{vlm_vmm_2022_emav}.

EMAV begon als een excel bestand.
Terwijl bij het Instituut voor Landbouw-, Visserij- en Voedingsonderzoek (ILVO) meer en meer data wordt verzameld,
begon de excel te groot te worden en niet meer onderhoudbaar.
Naast de grootte zijn er constant nieuwe wetenschappelijke inzichten die gevonden worden.
Hierdoor is een excel bestand niet genoeg meer.

Om dat probleem op te lossen is ILVO overgestapt naar een Windows Forms app, geschreven in C\#.
Dit is al een oud project en is al meerdere keren aangepast naar een betere, meer complexe versie.


De onderzoekers zijn vooral gekend in python. Ze hebben een basis van C\#, maar het is zeker niet makkelijk voor hun om de code van EMAV te navigeren.
Zij vinden problemen, foutjes in de resultaten, maar zij kunnen niet zomaar opzoeken waar het probleem zit.
Voorlopig maken ze Tickets aan via devops van Azure. Ze kunnen al veel info doorgeven via die beschrijving, maar het is nog steeds een zoektocht voor de it'ers om het probleem te vinden.

Zoals bij nieuwe features, tijdens het overbrengen van de info, gaat die communicatie niet altijd zo goed.
Vele onduidelijkheden vinden hun weg in de lijn van communicatie.

De documentatie wordt onderhouden door de onderzoekers, een word bestand met een paar flowcharts.
Zij weten dus niet exact hoe EMAV in elkaar zit.
Door die communicatie problemen kan het zijn dat wat de onderzoekers willen beschrijven in de documentatie mogelijks niet is hoe het is geimplementeerd.
Zeker bij het overstappen naar nieuwe versies van het model.
Hierdoor is er zeker sprake van Documentatie drift.


\section{Vaststelling}
